\section{From a frontier result to a banking capability}
\label{exec13:frontier-synthesis}

The inspected literature gives reasons to invest in better assistance, but it
does not yield a single ranking of causal systems. A system that extracts
relations from documents, a system that answers a supplied structural-model
question and a system that chooses a treatment policy may all receive a score
called causal accuracy. Their inputs, permissible assumptions and correct
outputs differ. A product proposal must recover these differences before it
uses an experimental result to select software or allocate development effort.

For this purpose, describe an evaluation claim by
\begin{equation}
 \mathcal E=(\mathcal T,\mathcal I,\mathcal A,\mathcal O,\mathcal B,\mathcal L),
 \label{exec13:frontier-claim}
\end{equation}
where $\mathcal T$ is the task distribution, $\mathcal I$ the information
supplied to the system, $\mathcal A$ its scientific assumptions, $\mathcal O$
the oracle or adjudication rule, $\mathcal B$ the tool, repair and human budget,
and $\mathcal L$ the loss used to score its output. This is a proposed
comparison vocabulary, not a theorem that these quantities are observed without
error. It asks a practical question: which parts of the published experiment
would remain the same in the proposed banking use? A change in any component
requires an argument or a new experiment before the original score can carry
the same interpretation.

For example, a graph provided in the prompt belongs to $\mathcal I$. Correctly
calculating a back-door formula from it tests conditional reasoning. Inferring
that graph from account and contact records is a different task with different
assumptions. Likewise, a human who repairs a mistaken join changes the combined
workflow's information and cost. The repaired answer can be valuable, but
attributing its correctness entirely to the model would misdescribe the product
being evaluated. This motivates an explicit comparison of the strongest
available alternatives rather than an agent-only leaderboard.

\subsection{Favorable and adverse results locate different opportunities}

The earliest broad frontier analysis treats language models as a channel for
domain knowledge and a possible aid to graph construction, critique and
robustness analysis. It also probes memorization and prompt sensitivity. These
observations motivate assistance with familiar concepts while warning that a
recognizable public dataset can supply information absent from a fresh bank
problem. Neither renaming a variable nor recalling a benchmark answer
establishes a new identification result \cite{v12_la_2305_00050}.

CLadder gives a more controlled formal task: generated questions from specified
causal structures with alternative verbalizations. That design isolates
important distinctions among association, intervention and counterfactual
reasoning. The supplied structure and generator, however, are part of the
experimental information. A bank workflow must also determine whether its
variable, treatment and population definitions support that structure. The
paper's remaining language-to-formal-tool problem is therefore central to the
proposed product, not an incidental user-interface issue
\cite{v12_la_2312_04350}.

Two comparisons show why the proposal should retain nonagent and smaller-model
baselines. In CausalEval, adding function calls improves the reported GPT-4o
CLADDER point score from 69 to 76 relative to ordinary chain-of-thought, but
reduces its e-CARE score from 90 to 73. These are outcomes from the named
historical configurations and samples. They demonstrate variation across tasks,
not a general tool penalty or benefit. In the inspected relation-extraction
study, the DDI F1 values for Single-Prompt C and PubMedBERT are 67.4 and 87.9.
The task concerns binary pairwise relations and the fine-tuned system receives
labeled training examples. The result supports considering supervised
extraction where the document domain and annotation budget are suitable; it
does not compare whole causal pipelines at equal total cost
\cite{exec13_census_causaleval,exec13_census_graph_validation}.

Retrieval has similarly mixed consequences. LACR combines background knowledge,
documents and optional numerical information. Its ASIA skeleton F1 falls from 1
to approximately .73 when documents are added, while other reported cases
benefit. Documents can add an applicable mechanism, introduce a conflicting
definition or repeat the same unsupported assertion. The scientific question
is which information changed and why it was relevant, rather than how many
documents entered the prompt \cite{v12_la_2402_15301}.

Auditing assumptions is another promising but restricted task. ARGUS assesses
the reported support for difference-in-differences assumptions, uses injected
flaws as local ground truth and explicitly avoids adjudicating causal truth.
Its five-paper human-label pilot and in-sample calibration do not establish
the error rate on a new collection of studies. A bank can use such an assistant
to find missing discussion or contradictory passages while keeping assignment,
parallel trends and measurement facts open to scientific assessment
\cite{v12_la_2609_30867}.

CausalCollab makes the longitudinal issue concrete. It defines an incremental
stylistic effect in sequential human--model interactions under explicit
positivity and sequential exchangeability premises. Its intervention mixes
text versions through a specified functional construction. A proposed
collections-message rewriting policy is not automatically that intervention:
it must define which actual messages are delivered, to whom, and after which
history. Model-generated counterfactual edits or outcome labels in an
evaluation remain generated evidence. The theorem's premises, the estimator's
approximation and the operational rewriting policy must be reviewed separately
\cite{exec13_census_collaboration}.

These results favor a portfolio of bounded roles. Template-driven formalization,
supervised extraction, retrieved evidence, formal graph algorithms and numerical
causal estimators can each be a baseline or component. A language model is
valuable where it reduces a measured bottleneck while preserving the meaning
of the task. Its presence is not a requirement for every component. The
architecture chapter develops how to enforce those boundaries in handoffs;
the question here is what evidence can justify selecting them.

\subsection{Task composition can reverse a published ranking}

Let $s_{ak}$ be method $a$'s success probability on task stratum $k$ and $w_k$
the proportion of eligible requests in that stratum, with $\sum_k w_k=1$.
Under a common success definition, the overall probability is
\begin{equation}
 S_a(w)=\sum_k w_k s_{ak}.
 \label{exec13:frontier-mixture}
\end{equation}
Even when every $s_{ak}$ is known, ranking methods requires $w$. A pooled score
from one benchmark mixture is not a mixture-free property of a system. When
the strata use incompatible outcomes, they should first remain separate rather
than be forced into this formula.

As an arithmetic illustration, use the two reported CausalEval point-score
pairs just discussed. Give CLADDER weight $w$ and e-CARE weight $1-w$. The
FunctionCall-minus-chain-of-thought difference is
\begin{equation}
 \Delta(w)=.07w-.17(1-w)=.24w-.17.
 \label{exec13:frontier-rank}
\end{equation}
It changes sign at $w=17/24\simeq.7083$. At equal weights the difference is
$-.05$; at $w=.8$ it is $.022$. This reweighting uses published point scores
solely to illustrate composition. It supplies neither their sampling
uncertainty nor a bank's task mix. The original paired outcomes and
mechanism-level dependence would be needed for an uncertainty calculation
\cite{exec13_census_causaleval,exec13_frontier_exact}.

The same principle applies to missing tasks. If all evaluated requests assume
an observed confounder set, the experiment cannot establish performance on
latent-confounding requests, whatever their future bank frequency. If the
support and weights of the intended workload are unknown, report stratum-level
results and plausible planning mixtures. A later public or bank intake sample
can estimate the mix. Changing the weighting after observing which method wins
would instead change the decision rule being evaluated.

The task family should be the meaningful sampling unit. Paraphrases from the
same mechanism, multiple prompt seeds and repeated tool calls share information.
The repeated-call variance example in the evaluation chapter explains why
their count cannot simply replace the count of independent mechanisms. For a
matched comparison, pair methods on the same held-out mechanisms, preserve all
retries and report uncertainty at the level where the evaluation law supplies
independence. No variance estimate can be recovered from a published average
alone.

\subsection{Component correctness must compose through conditional events}

Suppose a request passes five stages: correct formulation $E_1$, observation
construction $E_2$, identification $E_3$, numerical execution $E_4$ and
reporting $E_5$. By the chain rule,
\begin{equation}
 P\!\left(\bigcap_{j=1}^5 E_j\right)
 =P(E_1)\prod_{j=2}^5P\!\left(E_j\mid\bigcap_{k<j}E_k\right).
 \label{exec13:frontier-chain}
\end{equation}
This identity does not require independent stages. It does require that each
factor refer to the same workflow and condition on all relevant preceding
successes. Marginal accuracies reported in five unrelated papers are not those
conditional probabilities. They may concern different data, expert information
and definitions of success. Multiplying them would manufacture an end-to-end
claim from incompatible experiments.

For an exact planning example, conditional probabilities
$(.96,.95,.90,.97,.98)$ yield $.78025248$. If these numbers are instead only
marginal success probabilities in a common request population, the joint
success is bounded by
\begin{equation}
 \max\!\left(0,1-\sum_{j=1}^5P(E_j^c)\right)
 \leq P\!\left(\bigcap_{j=1}^5 E_j\right)
 \leq\min_jP(E_j),
 \label{exec13:frontier-frechet}
\end{equation}
giving $.76$ to $.90$. The lower bound follows from the union bound for
failures; the upper follows because joint success implies every individual
success. Two authored populations of 100 requests attain the endpoints. In
one, failure sets of sizes 4, 5, 10, 3 and 2 are disjoint, leaving 76 complete
successes. In the other the failure sets are nested inside the ten largest
failures, leaving 90. Both have exactly the same stage marginals
\cite{exec13_frontier_exact}.

For the proposal, this means stage tests diagnose where to improve, while a
separate composition test measures the completed request. A syntactically
valid graph can propagate a wrong population through every later stage. A
correct estimator may precisely calculate the wrong target. Conversely,
several stages can fail on the same small subset, so adding their failure
counts overstates the number of affected requests. Record per-request
transitions and preserve the joint history needed to distinguish these cases.

\subsection{An answer key is another measured object}

Let $Z$ be a benchmark's recorded label, $T$ the scientifically intended answer
and $\widehat T$ the system output. For exact-match correctness,
\begin{equation}
 \left|1\{\widehat T=T\}-1\{\widehat T=Z\}\right|
 \leq1\{T\ne Z\}.
 \label{exec13:frontier-oracle}
\end{equation}
Whenever the label is correct, both correctness indicators are equal;
otherwise their difference is at most one. Taking expectations therefore
bounds the difference between scientific accuracy and label agreement by the
answer-key error probability. If a defensible upper bound is $\epsilon$,
reported agreement $a$ implies only
$\max(0,a-\epsilon)\leq P(\widehat T=T)\leq\min(1,a+\epsilon)$.
This statement assumes a meaningful common answer space; set-valued and
ambiguous tasks require an explicitly defined acceptance relation.

With agreement $.9$ and a stipulated label error $.1$, truth accuracy can be
$.8$ or 1 in two simple 100-case constructions. Errors in the key can either
conceal wrong outputs or penalize correct ones. This is an exact example, not
an allegation that any inspected benchmark has a ten-percent error rate
\cite{exec13_frontier_exact}. The practical lesson is to audit the generator,
label rule and edge cases before using a score as product evidence. An
incorrect IV-to-population-ATE oracle cannot be repaired by making an agent
match that oracle more reliably. Likewise, model-generated labels and
model-generated answers can share errors even when their agreement is high.

For public observations without known causal truth, use the information that
is actually available: variable definitions, temporal admissibility, numerical
reproduction, sensitivity to assumptions and agreement with a justified design.
Do not silently promote a published observational point estimate into ground
truth. For generated mechanisms, keep generator truth outside the analysis
interface and verify the reference calculation independently. An intentional
nonidentified case must accept appropriate refusal or a valid set, rather than
reward a precise but unjustified number.

\subsection{Causal uncertainty matters through the decision}

Return to a bank choosing an action $a\in\{0,1\}$. Let $V(a,M)$ be its net
value under causal model $M$. When the evidence retains class $\mathcal M$,
one way to display the decision consequence is regret
\begin{equation}
 R(a,M)=\max_{b\in\{0,1\}}V(b,M)-V(a,M).
 \label{exec13:frontier-regret}
\end{equation}
A model's graph accuracy and its decision regret can order systems differently.
An edge irrelevant to the intervention query may contribute graph loss without
changing the decision; a single uncertain edge on the treatment-outcome path
can reverse the preferred action. A credible evaluation reports both the
scientific object and the loss appropriate to the decision.

The architecture chapter gives two Gaussian worlds with an identical observed
distribution but effects of 0 and $1/2$. Rescale the effect to dollars so that
the two candidate gains from acting are 0 and 500, and set action cost to 250.
Not acting has value zero in both worlds. Acting has value $-250$ or 250.
Because the available information has the same law in both worlds, any
randomized rule acts with the same marginal probability $p$. Its expected
regrets are $250p$ and $250(1-p)$, so
\begin{equation}
 \inf_{0\leq p\leq1}\max\{250p,250(1-p)\}=125.
 \label{exec13:frontier-minimax}
\end{equation}
Equality holds at $p=1/2$; either deterministic rule has worst regret 250.
This restricted two-world calculation illustrates the cost of unresolved
information. It is not a recommendation to randomize real banking actions,
and the world class and loss are authored. It explains why a request for an
informative experiment can be more valuable than another confident narrative
\cite{exec13_architecture_exact,exec13_frontier_exact}.

The next experiment should target information capable of distinguishing models
with different preferred actions. A public simulator can check that the
proposed procedure chooses such an experiment under known truth. Public
observations can establish feasibility of a data pipeline or stress an
assumption. Neither proves that the same model class, costs or experiment
availability holds at a particular bank. Those transport questions belong to
a later study justified by the prototype's findings.

\subsection{A public-evidence portfolio with explicit transfer limits}

An affordable portfolio follows the four banking journeys already developed.
For funding, public panel and macro records support data-clock, aggregation
and reproducibility checks, while declared generated laws supply causal truth
for uncertainty tests. For servicing, a controlled history/action generator
tests positivity failures, feedback, policy definitions and the effect of
withheld state information. For discovery, known mechanisms should include
observationally equivalent graphs whose effects differ, so a confident single
graph is not always rewarded. For observation construction, public or generated
documents can test extraction and linkage while preserving the difference
between a faithfully extracted assertion and a true causal statement.

Each journey should compare a deterministic reference, an established
statistical route and the proposed assistance where its role is relevant.
Give all arms the same admissible information and account separately for
training labels, retrieval, tool use and human repair. Freeze mechanisms,
splits, budgets and scoring before final evaluation. A fixed held-out set can
be used for that registered comparison; repeated development against its
answers creates a new dependence and requires a new final evaluation design.
An unavailable or malformed answer remains in the assigned-request denominator.

The interpretation should stay attached to the actual evidence. A document
review can justify an integration candidate. A local exact example can expose
a failure mechanism. A controlled prototype study can compare methods under
its declared law. A live-model study adds evidence about a particular model
revision and task distribution. A bank study adds evidence about the target
process and its decision consequences. These are useful, different claims;
their separation allows the proposal to advance now without purchasing an
expensive experiment merely to satisfy a deployment gate.

The current research frontier therefore supports a concrete proposal: build
typed formulation and evidence assistance around an auditable causal core;
retain supervised, deterministic and formal alternatives; measure composition
and review costs; and direct new information acquisition toward consequential
ambiguity. It leaves open the reliability of language-mediated assumptions,
performance outside benchmark support, annotation validity and the economics
of sustained human review. Those are falsifiable next questions. A count of
papers, agents or completed reports does not answer them.
\label{exec13:frontier-synthesis-end}
